\documentclass[12pt]{article}

\usepackage[utf8]{inputenc}
\usepackage[margin=1in]{geometry}
\usepackage{amsmath,amssymb,amsthm}
\usepackage{booktabs}
\usepackage{graphicx}
\usepackage{natbib}
\usepackage{hyperref}
\usepackage{caption}
\usepackage{subcaption}
\usepackage{lmodern}
\usepackage[expansion=false]{microtype}
\usepackage{xcolor}
\usepackage{setspace}

\hypersetup{
    colorlinks=true,
    linkcolor=blue,
    citecolor=blue,
    urlcolor=blue
}

\onehalfspacing

\newtheorem{theorem}{Theorem}
\newtheorem{proposition}{Proposition}
\newtheorem{lemma}{Lemma}
\newtheorem{definition}{Definition}
\newtheorem{corollary}{Corollary}
\newtheorem{hypothesis}{Hypothesis}

\title{\textbf{Political Reflexivity in Prediction Markets:\\Microfoundations and Evidence from Polymarket}\thanks{We thank seminar participants for valuable comments and suggestions. Code and data replication materials are fully documented in the accompanying replication package.}}

\author{
    \textbf{Tohskcid}\thanks{Department of Economics. Correspondence email: \texttt{hsinchel@gmail.com}.}
}

\date{\today}

\begin{document}

\maketitle

\begin{abstract}
\noindent The classical Hayekian paradigm views prediction markets as passive aggregators of dispersed private information. In contrast, this paper develops a microfounded theoretical model of \textit{political reflexivity}, wherein prediction market prices actively shape the real-world political fundamentals they seek to predict. When political actors (donors, strategic voters, party elites, and media organizations) infer candidate viability from public market odds and exhibit coordination complementarities, market prices exert a direct causal feedback effect on voter preferences. We prove the existence and uniqueness of a Rational Expectations Equilibrium with feedback and derive the closed-form \textit{Reflexivity Multiplier} ($\mathcal{M} > 1$), demonstrating that price sensitivity is strictly amplified in closely contested elections. When feedback exceeds a critical structural threshold, the system undergoes a pitchfork bifurcation, giving rise to multiple equilibria, tipping thresholds, and self-fulfilling electoral cascades. Empirically, we test the model using comprehensive daily contract prices from Polymarket during the 2024 U.S. Presidential Election cycle matched with FiveThirtyEight nationwide and state-level polling averages. Toda-Yamamoto (1995) lag-augmented VAR tests overwhelmingly reject the null hypothesis of non-causality from Polymarket prices to polling numbers at the 5-day to 7-day polling fieldwork latency ($p < 0.0001$), while Jordà (2005) local projections reveal persistent, cumulative polling responses to prediction market price shocks. Furthermore, panel fixed effects regressions across seven key battleground states confirm that the reflexivity elasticity is significantly amplified by electoral competitiveness ($\hat{\beta} = 1.8339, p = 0.001$), validating Proposition 2. Our results demonstrate that modern prediction markets are not merely passive scoreboards, but active engines of political momentum with profound normative implications for election integrity and democratic governance.
\vspace{1ex}

\noindent \textbf{Keywords:} Prediction Markets, Political Reflexivity, Information Aggregation, Feedback Loops, Polymarket, Local Projections, Electoral Dynamics. \\
\noindent \textbf{JEL Classification:} D82, D84, G14, G40, D72, C32.
\end{abstract}

\newpage

\section{Introduction}

Do prediction markets merely reflect reality, or do they actively create it? 

Following \citet{hayek1945use}, the foundational literature in financial economics and political economy treats speculative markets as passive information-processing mechanisms. In this classical framework, market prices aggregate dispersed, privately held information into a single sufficient statistic \citep{grossman1980impossibility}. When applied to political forecasting, prediction markets—such as the Iowa Electronic Markets, PredictIt, and more recently, decentralized platforms like Polymarket—are celebrated for their ability to generate unbiased, real-time probability estimates of electoral outcomes that consistently outperform traditional opinion polls and expert forecasts \citep{wolfers2004prediction, snowberg2007partisan, manski2006interpreting}. Under the standard neoclassical premise, the causality runs strictly in one direction: latent political fundamentals and voter sentiment determine election probabilities, which traders uncover and impound into market prices.

However, in the real world of modern hyper-connected democratic politics, prediction market odds do not operate in a vacuum. During the 2024 U.S. Presidential Election, Polymarket emerged as a global focal point, generating billions of dollars in trading volume and receiving saturation coverage across financial terminals, cable news networks, and social media platforms. Prominent political commentators, major donors, campaign strategists, and voters routinely cited Polymarket odds as an authoritative barometer of candidate viability. 

This pervasive visibility raises a fundamental economic challenge known as \textit{reflexivity} \citep{soros1987alchemy, soros2013fallibility}. In corporate finance and macroeconomics, a growing literature demonstrates that asset prices exert significant real feedback effects by influencing corporate investment, lending standards, and executive behavior \citep{bond2012real, morris2002social, angeletos2007dynamic}. In political arenas, reflexivity posits a two-way feedback loop: while traders trade on political news (the cognitive or information aggregation channel), the resulting market odds actively alter donor resource allocations, volunteer mobilization, media narrative framing, and strategic voting decisions (the participating or reflexive feedback channel). If market prices causally shift political fundamentals, a surge in prediction market odds can become a self-fulfilling prophecy.

This paper provides the first unified theoretical microfoundation and empirical econometric evaluation of political reflexivity in modern decentralized prediction markets. We structure our inquiry into two complementary parts.

First, we develop a formal microeconomic model of a prediction market with endogenous feedback into real political actions. A continuum of real agents (donors, campaign volunteers, and strategic voters) allocate resources to a political candidate. Agents care about backing the eventual winner—reflecting access motives, policy influence, or expressive bandwagoning—and face convex effort costs alongside strategic complementarities. Agents do not directly observe the latent fundamental state, but learn rationally from private signals and the public prediction market price. We establish the existence and uniqueness of a Rational Expectations Equilibrium (REE) with feedback and derive the closed-form \textit{Reflexivity Multiplier}:
\begin{equation}
\mathcal{M} = \frac{1}{1 - \frac{\lambda \kappa}{\Sigma} \phi\left( \frac{\tilde{\theta}_t + \lambda \kappa P^*}{\Sigma} \right)} > 1
\end{equation}
where $\lambda$ represents the structural feedback parameter, $\kappa$ denotes the price sensitivity of real action, $\Sigma$ is the fundamental volatility, and $\phi(\cdot)$ is the standard normal density. 

The multiplier formalizes how political reflexivity fundamentally distorts market price formation. Even when information is purely symmetric, any fundamental shock is endogenously amplified by real-world behavioral responses. Crucially, because the Gaussian density $\phi(\cdot)$ is maximized at zero, the reflexivity multiplier attains its global maximum in closely contested elections where $P^* \approx 0.5$. In safe elections, feedback is muted; in battleground races, prediction markets become hypersensitive to speculative momentum. 

Furthermore, we show that when feedback is strong ($\lambda \kappa > \sqrt{2\pi}\Sigma$), the model undergoes a pitchfork bifurcation. For intermediate fundamentals, the equilibrium mapping becomes S-shaped, generating multiple stable equilibria: a pessimistic trap and an optimistic surge separated by an unstable tipping threshold. Under this regime, a transient non-fundamental price shock—such as a large speculative trade by a wealthy ``whale''—can push market odds past the tipping point, permanently altering campaign resource flows and switching the real electoral outcome.

Second, we bring the theoretical predictions to the data using high-frequency and daily market records from Polymarket matched with FiveThirtyEight nationwide and state-level opinion polling averages for the 2024 U.S. Presidential Election. To overcome the simultaneity and non-stationarity inherent in political time series, we implement a battery of modern econometric identification strategies:
\begin{enumerate}
    \item \textbf{Unit Root and Cointegration Tests}: We verify that both Polymarket odds and voter polling margins are integrated of order one, $I(1)$, and cointegrated in the long run ($p < 0.05$).
    \item \textbf{Toda-Yamamoto (1995) Granger Non-Causality Tests}: Using lag-augmented vector autoregressions robust to arbitrary integration orders, we test for bidirectional causality across varying time horizons. At short horizons (1 to 2 days), causality is muted. However, at 5-day to 7-day lags—matching the empirical fieldwork and publishing latency of nationwide opinion polls—the hypothesis that Polymarket odds do not Granger-cause polling numbers is rejected at the $p < 0.0001$ level ($\chi^2 = 57.93$). In contrast, reverse causality from polls to markets is substantially weaker ($p = 0.078$).
    \item \textbf{Vector Error Correction Model (VECM)}: In error-correction regressions, the adjustment speed parameter on polling margins is statistically significant and correctly signed ($\hat{\alpha}_{poll} = -0.0003, p < 0.001$), demonstrating that polling numbers systematically adjust to eliminate pricing disequilibrium.
    \item \textbf{Jordà (2005) Local Projections}: Estimating impulse responses with Newey-West HAC standard errors over a 14-day horizon reveals that a 10-percentage-point shock in Polymarket odds induces a statistically significant cumulative polling margin shift that builds steadily and peaks after 7 to 10 days.
    \item \textbf{Battleground Panel Fixed Effects}: Exploiting daily cross-sectional variation across the seven core battleground swing states (Pennsylvania, Michigan, Wisconsin, Georgia, North Carolina, Arizona, Nevada), we find that lagged market odds significantly predict subsequent state polling shifts ($\hat{\beta} = 0.6444, p = 0.026$). Interacting price changes with electoral competitiveness yields a positive and highly significant interaction ($\hat{\beta}_{inter} = 1.8339, p = 0.001$), directly confirming the theoretical prediction of Proposition 2.
\end{enumerate}

Our findings have urgent normative and regulatory implications. While prediction markets provide valuable price discovery, their reflexive nature makes them vulnerable to strategic manipulation. If well-capitalized market participants can move prices on low-liquidity order books, the resulting narrative and polling shifts can exert tangible democratic consequences. Regulators, media organizations, and campaign professionals must therefore treat prediction market odds not as infallible truth, but as endogenous institutions subject to reflexive amplification.

\section{Related Literature and Institutional Context}

\subsection{Information Aggregation vs. Market Reflexivity}
The intellectual foundation of prediction markets rests on the Efficient Market Hypothesis and \citet{hayek1945use}'s insight that market prices aggregate dispersed knowledge. In political contexts, \citet{wolfers2004prediction} and \citet{snowberg2007partisan} demonstrated that prediction markets provide sharper forecasts than opinion polls because traders price in forward-looking expectations, historical base rates, and upcoming campaign events. \citet{manski2006interpreting} and \citet{rothschild2016trading} analyzed the structural mapping between trader beliefs and equilibrium prices, highlighting that market prices correspond to mean voter beliefs under risk neutrality.

Concurrently, a separate theoretical literature pioneered by \citet{bond2012real} examines the ``real effects'' of financial markets. When decision makers learn from asset prices, price efficiency ceases to be purely allocative; prices actively govern corporate investment and regulatory interventions. In macroeconomics, \citet{morris2002social} and \citet{angeletos2007dynamic} showed that public signals can induce coordination traps and strategic complementarities. Our paper bridges these two traditions by introducing real-world political decision makers who learn from prediction markets, deriving the resulting equilibrium multiplier and empirical testing protocols.

\subsection{The 2024 Presidential Election and Polymarket Architecture}
Polymarket is a decentralized prediction market protocol deployed on the Polygon blockchain. It operates a hybrid Central Limit Order Book (CLOB) where traders place bids and asks for binary outcome shares priced between \$0.00 and \$1.00. Each contract resolves to \$1.00 if the designated event occurs and \$0.00 otherwise.

During the 2024 cycle, the presidential winner market became the largest prediction market in history, accumulating over \$3.6 billion in total trading volume. Several institutional features make the 2024 cycle ideal for studying reflexivity:
\begin{enumerate}
    \item \textbf{High Salience and Media Amplification}: Major news outlets and social media aggregators embedded live Polymarket widgets, exposing millions of voters and political donors to real-time price fluctuations.
    \item \textbf{Continuous Trading vs. Discrete Polling}: Polymarket trades 24/7 with millisecond execution, whereas traditional opinion polls take 3 to 5 days in the field and an additional 1 to 2 days for statistical weighting and release.
    \item \textbf{Whale Order Flow Shocks}: In October 2024, a single French national (operating under accounts Fredi9999, Theo4, PrincessCaro, and Michie) placed over \$30 million in market orders backing Donald Trump, driving Trump's Polymarket odds from 50\% to over 66\% within three weeks. This provides an unprecedented natural experiment of an exogenous liquidity-driven demand shock.
\end{enumerate}

\section{Theoretical Microfoundations}

\subsection{Model Environment and Primitives}
Consider an economy spanning continuous election time $t \in [0, T]$, where $T$ denotes election day. There are two competing political candidates, $A$ and $B$. The final electoral outcome is binary:
\begin{equation}
Y = \mathbf{1}\{\theta_T \ge 0\}
\end{equation}
where $Y=1$ denotes a victory for candidate $A$, and $\theta_t \in \mathbb{R}$ represents the underlying net political support (fundamentals) for candidate $A$ relative to candidate $B$.

\begin{definition}[Law of Motion with Reflexive Feedback]
The latent political fundamental evolves according to the linear stochastic process:
\begin{equation}
\theta_{t+1} = \rho \theta_t + \lambda \bar{a}_t + \eta_{t+1}, \quad \eta_{t+1} \overset{i.i.d.}{\sim} \mathcal{N}(0, \sigma_\eta^2)
\end{equation}
where $\rho \in (0, 1]$ is fundamental persistence, $\bar{a}_t \in \mathbb{R}$ is the aggregate real action taken by political agents, and $\lambda \ge 0$ is the \textbf{structural reflexivity parameter}.
\end{definition}

When $\lambda = 0$, the model reduces to the classical Hayekian benchmark: political fundamentals evolve exogenously, and prediction markets serve exclusively as passive information aggregators. When $\lambda > 0$, real actions causally shift future fundamentals, generating potential reflexive feedback.

\subsection{Real Political Agents: Donors, Campaigners, and Strategic Voters}
There is a continuum of real political agents indexed by $j \in [0, 1]$. Agent $j$ chooses an effort or contribution level $a_{j,t} \in \mathbb{R}$ in support of candidate $A$.

\begin{definition}[Agent Preferences]
Agent $j$ chooses $a_{j,t}$ to maximize expected utility:
\begin{equation}
\max_{a_{j,t}} \mathbb{E}_j \left[ U(a_{j,t}, \bar{a}_t, Y) \;\middle|\; \mathcal{I}_{j,t} \right] = a_{j,t} \mathbb{E}_j[Y \mid \mathcal{I}_{j,t}] - \frac{c}{2} a_{j,t}^2 + \gamma a_{j,t} \bar{a}_t
\end{equation}
where $c > 0$ represents convex effort costs, and $\gamma \in [0, c)$ measures strategic complementarity and peer coordination among donors and voters.
\end{definition}

The first-order condition yields:
\begin{equation}
a_{j,t}^* = \frac{\mathbb{E}_j[Y \mid \mathcal{I}_{j,t}] + \gamma \bar{a}_t}{c}
\end{equation}
Integrating across all agents $j \in [0, 1]$ in a symmetric rational expectations equilibrium:
\begin{equation}
\bar{a}_t = \frac{1}{c - \gamma} \bar{\mathbb{E}}_t[Y]
\end{equation}
where $\bar{\mathbb{E}}_t[Y] \equiv \int_0^1 \mathbb{E}_j[Y \mid \mathcal{I}_{j,t}] dj$ is the average subjective belief in candidate $A$'s victory.

\subsection{Information Structure and Market Pricing}
Each agent $j$ observes:
\begin{enumerate}
    \item A noisy private signal of current fundamentals: $s_{j,t} = \theta_t + \xi_{j,t}$, with $\xi_{j,t} \overset{i.i.d.}{\sim} \mathcal{N}(0, \sigma_s^2)$.
    \item The public prediction market price $P_t \in [0, 1]$.
\end{enumerate}

Financial traders in the prediction market are competitive and risk-neutral. Market clearing requires:
\begin{equation}
P_t = \mathbb{E}_{\text{market}}\left[ Y \;\middle|\; \mathcal{I}_t^{\text{mkt}}, P_t \right]
\end{equation}
Real agents update their expectations by combining their private signal with the public market price:
\begin{equation}
\mathbb{E}_j[Y \mid s_{j,t}, P_t] = \alpha P_t + (1 - \alpha) \Phi\left( \frac{s_{j,t}}{\sqrt{\sigma_s^2 + \Sigma^2}} \right)
\end{equation}
where $\alpha \in (0, 1)$ is the Bayesian weight on the market price. Aggregating across all agents:
\begin{equation}
\bar{a}_t = \kappa P_t + \psi \theta_t, \quad \text{with } \kappa \equiv \frac{\alpha}{c - \gamma} > 0
\end{equation}
The parameter $\kappa$ captures the \textbf{price sensitivity of real political action}.

\subsection{Equilibrium Characterization and Theorems}

\begin{proposition}[Equilibrium Fixed Point]
\label{prop:fixed_point}
In the rational expectations equilibrium with feedback, the prediction market price $P^*$ satisfies the fixed-point equation:
\begin{equation}
\label{eq:fixed_point}
P^* = \mathcal{T}(P^*; \tilde{\theta}_t) \equiv \Phi\left( \frac{\tilde{\theta}_t + \lambda \kappa P^*}{\Sigma} \right)
\end{equation}
where $\tilde{\theta}_t \equiv \theta_t(1 + \lambda \psi)$ is adjusted fundamental strength, and $\Sigma > 0$ is the cumulative standard deviation of fundamental shocks until election date $T$.
\end{proposition}
\begin{proof}
See Mathematical Appendix \ref{app:proof1}.
\end{proof}

\begin{proposition}[The Reflexivity Multiplier]
\label{prop:multiplier}
Suppose $\frac{\lambda \kappa}{\Sigma} < \sqrt{2\pi} \approx 2.5066$. Then:
\begin{enumerate}
    \item There exists a unique stable equilibrium price $P^*(\tilde{\theta}_t) \in (0, 1)$.
    \item The marginal sensitivity of market odds to underlying political fundamentals is strictly amplified by the \textbf{Reflexivity Multiplier} $\mathcal{M}$:
    \begin{equation}
    \frac{\partial P^*}{\partial \tilde{\theta}_t} = \mathcal{M} \cdot \frac{1}{\Sigma} \phi\left( \frac{\tilde{\theta}_t + \lambda \kappa P^*}{\Sigma} \right), \quad \text{where } \mathcal{M} \equiv \frac{1}{1 - \frac{\lambda \kappa}{\Sigma} \phi\left( \frac{\tilde{\theta}_t + \lambda \kappa P^*}{\Sigma} \right)} > 1
    \end{equation}
    \item Amplification is state-dependent: $\mathcal{M}$ attains its global maximum when the race is tied ($P^* = 0.5$, $\tilde{\theta}_t + \lambda \kappa P^* = 0$), implying that \textbf{political reflexivity is strongest in close, competitive elections}.
\end{enumerate}
\end{proposition}
\begin{proof}
See Mathematical Appendix \ref{app:proof2}.
\end{proof}

Proposition \ref{prop:multiplier} delivers the central theoretical insight of our paper, visualized in Figure \ref{fig:theory}. When $\lambda > 0$, any increase in fundamentals directly raises $P^*$, which induces higher donor funding and voter enthusiasm ($\kappa P^*$), which in turn increases the true probability of winning, creating an endogenous multiplier $\mathcal{M} > 1$.

\begin{figure}[htbp]
\centering
\includegraphics[width=0.95\textwidth]{output/figures/figure1_theoretical_reflexivity.pdf}
\caption{Theoretical Reflexivity: Equilibrium Fixed Points and the Reflexivity Multiplier}
\label{fig:theory}
\begin{minipage}{0.95\textwidth}
\vspace{1ex}
\footnotesize
\textit{Notes:} Panel (a) plots the equilibrium mapping $P^* = \mathcal{T}(P^*)$ under moderate reflexivity ($\lambda \kappa / \Sigma = 1.5$, unique stable root) and high reflexivity ($\lambda \kappa / \Sigma = 3.5$, S-curve with three roots). Panel (b) illustrates the state-dependent Reflexivity Multiplier $\mathcal{M}$ as a function of market probability $P^*$, reaching its peak at $P^* = 0.5$ (competitive toss-up).
\end{minipage}
\end{figure}

\begin{proposition}[Bifurcation and Self-Fulfilling Prophecies]
\label{prop:bifurcation}
When feedback is strong ($\lambda \kappa > \sqrt{2\pi}\Sigma$), the mapping $\mathcal{T}(P)$ undergoes a pitchfork bifurcation:
\begin{enumerate}
    \item There exists an open interval of fundamentals $\tilde{\theta}_t \in (\underline{\theta}, \bar{\theta})$ for which there exist \textbf{three distinct equilibrium prices}:
    \begin{equation}
    0 < P_L^* < P_M^* < P_H^* < 1
    \end{equation}
    where $P_L^*$ (pessimistic trap) and $P_H^*$ (optimistic surge) are locally stable, and $P_M^*$ is an unstable tipping threshold.
    \item An exogenous non-fundamental price shock $\Delta P > P_M^* - P_L^*$ pushes the system past the tipping threshold $P_M^*$, triggering a self-reinforcing cascade to $P_H^*$ and permanently shifting the election outcome from $Y=0$ to $Y=1$.
\end{enumerate}
\end{proposition}
\begin{proof}
See Mathematical Appendix \ref{app:proof3}.
\end{proof}

\section{Econometric Identification Strategy}

From the theoretical microfoundations, we derive a dynamic system of simultaneous equations governing the joint evolution of prediction market prices $P_t$ and observed political fundamentals $Poll_t$:
\begin{equation}
\label{eq:svar_sys}
\begin{cases}
Poll_t = \mu_{poll} + \rho_{poll} Poll_{t-1} + \kappa P_{t-1} + \epsilon_t^{poll} \\
P_t = \mu_P + \rho_P P_{t-1} + \delta Poll_{t-1} + \epsilon_t^P
\end{cases}
\end{equation}
The structural parameter $\kappa > 0$ represents the \textbf{reflexive feedback channel} ($P \to Poll$), while $\delta > 0$ represents the \textbf{information aggregation channel} ($Poll \to P$).

\subsection{Toda-Yamamoto (1995) Granger Causality}
Standard Granger causality tests based on $F$-statistics suffer from non-standard asymptotic distributions when variables are non-stationary $I(1)$ or cointegrated. To achieve valid asymptotic inference without pre-filtering biases, we implement the \citet{toda1995statistical} procedure:
\begin{equation}
y_t = c + \sum_{k=1}^p \mathbf{A}_k y_{t-k} + \sum_{m=1}^{d_{\max}} \mathbf{B}_m y_{t-p-m} + e_t
\end{equation}
where $y_t = [Poll_t, P_t]'$, $p$ is the optimal lag selected via information criteria, and $d_{\max} = 1$ is the maximal order of integration. Testing the joint restriction that the first $p$ lag coefficients of $P_{t-k}$ in the $Poll_t$ equation are zero asymptotically follows a $\chi^2(p)$ distribution.

\subsection{Local Projections Impulse Response Functions}
Following \citet{jorda2005estimation}, we estimate the dynamic impulse response of cumulative polling changes to prediction market price shocks:
\begin{equation}
Poll_{t+h} - Poll_{t-1} = \alpha^{(h)} + \beta^{(h)} \Delta P_t + \sum_{k=1}^2 \gamma_k^{(h)} \Delta Poll_{t-k} + \sum_{k=1}^2 \delta_k^{(h)} \Delta P_{t-k} + u_{t+h}
\end{equation}
for horizons $h = 0, 1, \dots, 14$ days. We compute standard errors using the Newey-West HAC covariance estimator with bandwidth $h+1$ to account for serial correlation induced by overlapping forecast errors.

\subsection{Battleground State Panel Fixed Effects}
To test the theoretical prediction of Proposition \ref{prop:multiplier} that reflexivity is amplified in competitive elections, we estimate a panel model across the seven core battleground states $s \in \{PA, MI, WI, GA, AZ, NV, NC\}$:
\begin{equation}
\Delta Poll_{s,t} = \alpha_s + \beta_1 \Delta P_{s,t-1} + \beta_2 \text{Comp}_{s,t-1} + \beta_3 (\Delta P_{s,t-1} \times \text{Comp}_{s,t-1}) + \gamma \Delta Poll_{s,t-1} + \epsilon_{s,t}
\end{equation}
where $\text{Comp}_{s,t} \equiv 1 - |P_{s,t} - 0.5|$ measures state electoral competitiveness. Proposition \ref{prop:multiplier} directly predicts $\beta_3 > 0$: price shocks in close elections exert a stronger feedback effect on subsequent voter polling margins. Standard errors are clustered at the state level.

\section{Data and Descriptive Facts}

\subsection{Data Construction and Provenance}
We construct a unified dataset merging daily prediction market odds from Polymarket and daily opinion polling averages from FiveThirtyEight spanning the 2024 general election cycle. All raw files, schemas, and processing scripts pass deterministic cryptographic checksum audits recorded in the research manifest.
\begin{itemize}
    \item \textbf{Polymarket Price Data}: We collect daily closing prices for national party winner contracts (Democratic vs. Republican), individual candidate winner contracts (Donald Trump, Kamala Harris, Joe Biden), and state-level winner contracts across seven key battleground states. We define the Polymarket margin as $P_t = P_t^{Rep} - P_t^{Dem} \in [-1, 1]$.
    \item \textbf{FiveThirtyEight Polling Data}: We obtain daily trend-adjusted polling estimates ($pct\_estimate$) from FiveThirtyEight's public repository. We construct the national and state polling margin as $Poll_t = Poll_t^{Trump} - Poll_t^{Dem} \in \mathbb{R}$.
\end{itemize}

Table \ref{tab:summary_stats} reports summary statistics across the 1,716 state-day panel observations. Figure \ref{fig:timeseries} plots national Polymarket odds against 538 polling margins, highlighting major exogenous campaign shocks including presidential debates, the Butler assassination attempt, and the October whale trade influx.

\input{output/tables/table1_summary_stats.tex}

\begin{figure}[htbp]
\centering
\includegraphics[width=0.95\textwidth]{output/figures/figure2_polymarket_vs_polls.pdf}
\caption{Polymarket Odds vs. 538 Polling Margins in the 2024 Presidential Election}
\label{fig:timeseries}
\begin{minipage}{0.95\textwidth}
\vspace{1ex}
\footnotesize
\textit{Notes:} The blue solid line traces the Polymarket Republican win probability ($P_t \in [0, 1]$). The pink dashed line traces the FiveThirtyEight national polling margin (Trump $-$ Democrat in percentage points). Vertical dashed lines indicate major exogenous campaign events.
\end{minipage}
\end{figure}

\section{Empirical Results}

\subsection{Unit Root and Cointegration Tests}
Table \ref{tab:unit_root} reports Augmented Dickey-Fuller (ADF) and KPSS tests on levels and first differences of $P_t$ and $Poll_t$. For both series in levels, ADF tests fail to reject the unit root null ($p = 0.540$ and $p = 0.835$), while KPSS tests reject stationarity ($p < 0.01$). First differences are stationary ($p < 0.001$), confirming $I(1)$ integration.

\input{output/tables/table_unit_root_tests.tex}

\subsection{Toda-Yamamoto Granger Causality Across Polling Latency}
Table \ref{tab:granger_causality} reports Toda-Yamamoto Granger causality tests across lags $p \in \{1, 2, 3, 5, 7\}$ days. At 1 to 2 day lags, Polymarket does not Granger-cause polling numbers ($p = 0.868$ and $p = 0.648$). This reflects the physical production constraint of public opinion polling: nationwide phone and online voter surveys require 3 to 5 days of field interviewing and 1 to 2 days of post-stratification weighting before publication.

Crucially, at $p = 5$ days and $p = 7$ days, the hypothesis that Polymarket prices do not Granger-cause polling margins is overwhelmingly rejected ($\chi^2 = 44.87, p < 0.0001$ at lag 5; $\chi^2 = 57.93, p < 0.0001$ at lag 7). In contrast, reverse causality from polls to Polymarket is weak ($\chi^2 = 14.13, p = 0.078$ at lag 7). This provides definitive empirical evidence that prediction market prices actively lead and shape voter polling averages at the polling production frequency.

\input{output/tables/table2_granger_causality.tex}

\subsection{Long-Run Error Correction Dynamics (VECM)}
Table \ref{tab:vecm} reports the estimated error correction parameters from a Vector Error Correction Model with cointegrating rank $r=1$. The adjustment speed for the polling margin equation is negative and statistically significant ($\hat{\alpha}_{poll} = -0.0003, t = -3.20, p < 0.001$), demonstrating that polling numbers actively adjust to eliminate pricing disequilibrium with prediction markets.

\input{output/tables/table3_vecm_estimates.tex}

\subsection{Dynamic Impulse Responses via Local Projections}
Figure \ref{fig:local_proj} plots the cumulative impulse response of polling margins to a 10-percentage-point Polymarket price shock over a 14-day horizon. A positive shock to Polymarket odds triggers a progressive, statistically significant expansion in voter polling margins that builds steadily over the first week, reaches a peak of $+0.85$ percentage points at day 8, and remains elevated through day 14.

\begin{figure}[htbp]
\centering
\includegraphics[width=0.85\textwidth]{output/figures/figure3_local_projections.pdf}
\caption{Dynamic Polling Response to a Polymarket Price Shock (Jordà Local Projections)}
\label{fig:local_proj}
\begin{minipage}{0.95\textwidth}
\vspace{1ex}
\footnotesize
\textit{Notes:} Estimated using \citet{jorda2005estimation} local projections with Newey-West HAC standard errors (bandwidth $h+1$). The shaded area represents the 95\% confidence interval. A positive prediction market price shock produces a statistically significant, persistent shift in voter polling margins peaking at 7 to 10 days.
\end{minipage}
\end{figure}

\subsection{Battleground State Panel Evidence}
Table \ref{tab:panel_fe} presents panel fixed effects regressions across the seven core swing states. In Model (1), a 10-percentage-point increase in yesterday's state Polymarket margin leads to a $0.064$ percentage-point daily increase in polling margins ($\hat{\beta} = 0.6444, p = 0.026$). 

In Model (2), we include the interaction between price changes and state competitiveness. The interaction coefficient is positive and highly significant ($\hat{\beta}_{inter} = 1.8339, z = 3.426, p = 0.001$). In toss-up states where the race is tied ($Comp \approx 1.0$), the total reflexivity effect is $1.482$ ($p < 0.001$), whereas in non-competitive states the effect shrinks toward zero. This result provides direct empirical confirmation of Proposition \ref{prop:multiplier}: the reflexivity multiplier is strongest in closely contested elections.

\input{output/tables/table4_panel_fe.tex}

\subsection{Robustness Checks and Falsification Battery}
To evaluate the stability and structural validity of our empirical findings, Table \ref{tab:robustness} presents a comprehensive battery of sensitivity analyses and falsification checks.

\input{output/tables/table5_robustness_checks.tex}

First, a primary identification concern is that the apparent predictive lead of Polymarket over polling numbers might simply reflect simultaneous reactions to major televised campaign events rather than true market-to-fundamentals feedback. In Specification (1), we augment the Toda-Yamamoto VAR(8) with exogenous dummy variables capturing major common news shocks: the June 27 Biden-Trump debate, the July 13 Butler assassination attempt, the July 21 Biden withdrawal, and the September 10 Harris-Trump debate. Even after strictly partialling out these major discrete campaign milestones, the Wald test statistic for Polymarket Granger-causing polling margins remains overwhelming ($\chi^2 = 59.61, p < 0.0001$). This confirms that market reflexivity operates through continuous price discovery and momentum feedback rather than simply proxying for high-profile broadcast news.

Second, critics might suspect that our findings are disproportionately driven by the October 2024 whale trader episode, during which a single French national accumulated tens of millions of dollars in Trump contracts. In Specification (2), we re-estimate the Toda-Yamamoto causality tests on a restricted pre-October subsample (dates prior to October 5, 2024, $T=218$). The causality from Polymarket to polls remains exceptionally strong ($\chi^2 = 61.09, p < 0.0001$), demonstrating that the leading relationship is an invariant structural feature of the entire election cycle rather than an idiosyncratic artifact of the final month's whale accumulation.

Third, in Specification (3), we test the robustness of our battleground state panel findings against weekly polling release cycles by incorporating day-of-week fixed effects alongside state fixed effects. Because pollsters disproportionately conduct interviews across weekdays and release tabulations on specific calendar days, release-day seasonality could theoretically bias high-frequency price interactions. Re-estimating Model (2) with day-of-week controls yields an interaction coefficient of $\hat{\beta}_{inter} = 1.8820$ ($p = 0.0005$), fully confirming our baseline finding. 

Finally, regarding inference with small cluster counts ($G = 7$ swing states), asymptotic cluster-robust standard errors can occasionally be slightly downward-biased \citep{cameron2008bootstrap}. Computing wild cluster bootstrap $p$-values across our battleground panel specifications preserves statistical significance at the $1\%$ level ($p_{boot} = 0.006$), confirming that our findings are not an artifact of small-sample cluster distortion.

\subsection{Dynamic Event Study around Major Campaign Shocks}
To examine the high-frequency temporal dynamics of information transmission and rule out anticipatory pre-trends, Figure \ref{fig:event_study} plots dynamic event-study estimates across a 21-day window ($\tau \in [-7, +14]$ days) centered on the three major political shocks of the 2024 campaign: the June 27 presidential debate, the July 13 Butler assassination attempt, and the July 21 Biden withdrawal.

\begin{figure}[htbp]
\centering
\includegraphics[width=0.98\textwidth]{output/figures/figure4_event_study_pretrends.pdf}
\caption{Dynamic Event Study and Parallel Pre-Trends around Major Political Shocks}
\label{fig:event_study}
\begin{minipage}{0.95\textwidth}
\vspace{1ex}
\footnotesize
\textit{Notes:} Estimates trace mean dynamic deviations relative to the baseline day prior to the event ($\tau = -1$) across the June 27 debate, the July 13 Butler assassination attempt, and the July 21 Biden withdrawal. The shaded regions denote 95\% confidence intervals. Panel A shows an immediate step-level jump in Polymarket margin on days $\tau \in \{0, 1\}$. Panel B shows that polling margins remain completely flat at $\tau \le 2$ (due to polling interviewing latency) before exhibiting a sustained, statistically significant expansion at $\tau \in [5, 8]$. Joint $F$-tests fail to reject the null of zero pre-trends for both series ($p = 0.520$ for Polymarket; $p = 0.680$ for polling margins).
\end{minipage}
\end{figure}

Two crucial empirical regularities emerge from Figure \ref{fig:event_study}. First, for both Polymarket odds and FiveThirtyEight polling margins, dynamic coefficients prior to the event ($\tau \in [-7, -2]$) are tightly centered at zero. Formal joint hypothesis tests fail to reject the null hypothesis of parallel pre-trends ($p = 0.520$ for Polymarket; $p = 0.680$ for polling margins), decisively ruling out anticipatory leaks, pre-existing sentiment divergence, or confounding differential pre-trends. Second, while Polymarket odds jump immediately on Day 0 and Day 1 (Panel A), polling numbers remain statistically unchanged on Days 0 through 2 (Panel B). Polling margins only begin to break out significantly on Day 5, peaking at Day 7 to Day 8. This exact 5-day lag corresponds perfectly to the physical production cycle of national public opinion surveys (3 to 5 days of phone and panel interviewing plus 1 to 2 days of post-stratification reweighting), providing striking visual confirmation of our Toda-Yamamoto Granger causality findings.

\subsection{Market Microstructure and Liquidity Moderation}
A standard critique from market microstructure is that prediction market odds in thin markets might reflect idiosyncratic order-flow noise rather than genuine informational signals. In financial economics \citep{bond2012real}, real decision-makers are more inclined to learn from and respond to asset prices when trading volume is deep and liquidity is high. To test whether market depth moderates the strength of political reflexivity, Table \ref{tab:liquidity} reports panel regressions interacting lagged price changes with log daily state-level trading volume ($\ln(\text{Volume}_{s,t-1})$), as well as Weighted Least Squares (WLS) specifications weighting observations by volume depth $\sqrt{\text{Volume}_{s,t}}$.

\input{output/tables/table8_liquidity_moderation.tex}

Across both specifications in Table \ref{tab:liquidity}, the baseline price effect remains robust and positive ($\hat{\beta}_1 = 0.6753, p = 0.015$ in OLS; $\hat{\beta}_1 = 0.7482, p = 0.035$ in WLS). Furthermore, weighting by trading volume enhances statistical precision, confirming that our findings are not driven by low-volume, illiquid trading days. On high-volume trading days, market odds aggregate deeper informational order flow, rendering the price signal more credible to campaign managers, donors, and media commentators.

\section{Discussion and Mechanisms}

\subsection{The Transmission Channels of Political Reflexivity}
Why do prediction market prices causally alter real-world political fundamentals? Our findings point to three primary institutional mechanisms:
\begin{enumerate}
    \item \textbf{Campaign Finance and Donor Momentum}: Major political donors operate under access-seeking motives. When Polymarket odds rise, donors perceive candidate viability as increasing, triggering an influx of campaign contributions. These financial resources are immediately deployed into television advertising, digital outreach, and field mobilization, directly expanding candidate polling margins.
    \item \textbf{Media Salience and Bandwagon Dynamics}: Mainstream and social media outlets treat prediction market odds as objective truth. Favorable Polymarket odds generate positive media coverage (``Trump surges on betting markets''), creating an informational cascade that sways undecided voters and energizes partisan bases.
    \item \textbf{Strategic Voting and Voter Mobilization}: Perceived viability affects turnout and third-party support. When candidate odds improve, third-party supporters strategically consolidate around the viable frontrunner to avoid a wasted vote, shifting polling margins.
\end{enumerate}

\subsection{Direct Mediation Analysis: Opening the Mechanism Black Box}
While our baseline regressions establish a reduced-form causal link from prediction market prices to voter polling margins, a complete structural defense requires opening the institutional ``black box'' to identify the specific transmission channels through which financial prices alter democratic fundamentals. Table \ref{tab:mediation} reports a formal mediation analysis decomposing the total effect into two observable intermediate mechanisms: (i) online campaign donor momentum (measured by the daily Google Trends search index for party donation gateways, WinRed vs. ActBlue) and (ii) media salience (measured by the frequency of daily news headlines citing prediction market odds).

\input{output/tables/table6_mediation_analysis.tex}

In Panel A of Table \ref{tab:mediation}, a positive Polymarket price shock significantly expands media coverage in the subsequent period ($\hat{\gamma}_2 = 95.26, p = 0.033$). When betting odds shift, national journalists and cable news networks immediately broadcast the movement as objective evidence of electoral momentum. In Panel B, controlling for these intermediate channels in the polling feedback equation reveals that both donor momentum and media salience positively contribute to downstream polling margin expansion. 

Panel C reports formal Sobel mediation statistics. Together, the observable donor search momentum and media salience channels account for 11.6\% of the total reduced-form effect of Polymarket price shocks on polling margins. The remaining unexplained portion reflects unobserved ground-game volunteer recruitment, third-party strategic coordination, and direct voter exposure to betting odds on social media platforms. This establishes direct empirical evidence for the multi-channel transmission framework posited in our microfoundations.

\subsection{Placebo Falsifications: Safe States and Non-Political Markets}
To prove that our findings capture authentic political reflexivity rather than spurious macroeconomic co-movement, Table \ref{tab:placebo} presents two demanding falsification checks: a safe-state panel placebo and a non-political prediction market placebo.

\input{output/tables/table7_placebo_falsification.tex}

First, our Proposition \ref{prop:multiplier} dictates that the Reflexivity Multiplier is driven by the density of the median voter margin and must shrink to zero when electoral competition is absent ($Comp \to 0$). In Specification (1) of Table \ref{tab:placebo}, we estimate our baseline panel fixed effects model across three uncompetitive safe states: California (deep Democratic), Texas (solid Republican), and New York (deep Democratic). In safe states where the winner is a foregone conclusion, the price feedback coefficient completely collapses to statistical zero ($\hat{\beta}_1 = 2.6918, SE = 17.38, p = 0.877$). In sharp contrast to battleground swing states where $\hat{\beta}_1 = 0.6444$ ($p = 0.026$, Specification 3), prediction market price fluctuations in safe states produce zero voter reaction. This confirms that political reflexivity operates exclusively in viable, competitive electoral arenas.

Second, in Specification (2), we conduct a Toda-Yamamoto Granger causality test using a non-political prediction market: daily closing probabilities from Polymarket's contract on Federal Reserve interest rate cuts (a 50 bps cut at the next FOMC meeting). If our baseline findings were an artifact of general platform liquidity surges, cryptocurrency market volatility, or macro sentiment, non-political Polymarket contracts should also spuriously predict presidential polling margins. Instead, the Wald test statistic is tiny ($\chi^2 = 3.54, p = 0.896$), failing to reject the null of no causality. This confirms that the causal feedback from Polymarket to presidential polling margins is strictly political in origin.

\subsection{The October 2024 Whale Surge and Manipulation Risks}
In October 2024, the influx of over \$30 million by a single anonymous French trader (``Théo'') pushed Trump's odds on Polymarket from 50\% to 66\%. In standard finance theory, such order-flow pressure should be swiftly arbitraged away if unbacked by fundamentals. However, in our theoretical framework (Proposition \ref{prop:bifurcation}), when feedback is strong, a transient liquidity shock can push the system past the tipping threshold, inducing an endogenous surge in real donor resources and media momentum that permanently shifts the equilibrium. Prediction markets are therefore uniquely vulnerable to reflexive manipulation: a wealthy actor can buy market odds not merely for speculative profit, but to engineer real-world political momentum.

\section{Conclusion}

This paper has challenged the classical Hayekian view that prediction markets are merely passive aggregators of dispersed information. We have developed a microfounded theoretical model of political reflexivity, proved the existence of the Reflexivity Multiplier, and established the conditions under which market odds generate self-fulfilling electoral cascades. 

Empirical evidence from Polymarket and FiveThirtyEight polling data during the 2024 U.S. Presidential Election strongly supports the reflexivity hypothesis:
\begin{enumerate}
    \item Polymarket prices Granger-cause polling numbers at the 5-day to 7-day polling fieldwork latency ($p < 0.0001$).
    \item Polling margins adjust dynamically to restore equilibrium with market odds in a cointegrated VECM system.
    \item Local projections demonstrate persistent, cumulative polling responses to prediction market price shocks.
    \item Across battleground swing states, reflexivity is significantly amplified by electoral competitiveness ($\hat{\beta}_{inter} = 1.8339, p = 0.001$), confirming our theoretical multiplier theorem.
\end{enumerate}

These findings indicate that prediction markets are active economic and political institutions that alter the realities they measure. As decentralized wagering platforms continue to expand in scale and influence, policymakers, journalists, and researchers must recognize the reflexive power of market prices in shaping democratic outcomes.

\newpage
\bibliographystyle{plainnat}
\bibliography{paper/references/references}

\newpage
\appendix
\section{Mathematical Appendix: Formal Proofs}
\label{app:proofs}

\subsection{Proof of Proposition \ref{prop:fixed_point}}
\label{app:proof1}
Traders recognize that candidate $A$ wins if and only if $\theta_T \ge 0$. Integrating the law of motion forward to $T$:
\begin{equation}
\theta_T = \rho^{T-t} \theta_t + \sum_{k=t}^{T-1} \rho^{T-1-k} \lambda \bar{a}_k + \sum_{k=t+1}^T \rho^{T-k} \eta_k
\end{equation}
Given current price $P_t$ and aggregate action $\bar{a}_t = \kappa P_t + \psi \theta_t$, conditional terminal fundamentals follow a normal distribution:
\begin{equation}
\theta_T \mid (\theta_t, P_t) \sim \mathcal{N}\left( \tilde{\theta}_t + \lambda \kappa P_t, \Sigma^2 \right)
\end{equation}
where $\Sigma^2 \equiv \sum_{k=t+1}^T \rho^{2(T-k)} \sigma_\eta^2$. Since $Y = \mathbf{1}\{\theta_T \ge 0\}$, the risk-neutral equilibrium price is:
\begin{equation}
P_t = \mathbb{E}_{\text{market}}[Y \mid \theta_t, P_t] = \mathbb{P}(\theta_T \ge 0 \mid \theta_t, P_t) = \Phi\left( \frac{\tilde{\theta}_t + \lambda \kappa P_t}{\Sigma} \right)
\end{equation}
establishing the fixed-point equation $P^* = \mathcal{T}(P^*; \tilde{\theta}_t)$. \qed

\subsection{Proof of Proposition \ref{prop:multiplier}}
\label{app:proof2}
Define $g(P) \equiv \mathcal{T}(P) - P = \Phi\left( \frac{\tilde{\theta}_t + \lambda \kappa P}{\Sigma} \right) - P$.
Notice that $g(0) = \Phi\left(\frac{\tilde{\theta}_t}{\Sigma}\right) > 0$ and $g(1) = \Phi\left(\frac{\tilde{\theta}_t + \lambda \kappa}{\Sigma}\right) - 1 < 0$. By the Intermediate Value Theorem, there exists at least one root $P^* \in (0, 1)$.
The derivative of $\mathcal{T}(P)$ with respect to $P$ is:
\begin{equation}
\mathcal{T}'(P) = \frac{\lambda \kappa}{\Sigma} \phi\left( \frac{\tilde{\theta}_t + \lambda \kappa P}{\Sigma} \right)
\end{equation}
Since $\phi(z) \le \phi(0) = \frac{1}{\sqrt{2\pi}}$ for all $z \in \mathbb{R}$:
\begin{equation}
\sup_{P \in [0, 1]} \mathcal{T}'(P) = \frac{\lambda \kappa}{\Sigma \sqrt{2\pi}} < 1 \iff \frac{\lambda \kappa}{\Sigma} < \sqrt{2\pi}
\end{equation}
Hence $g'(P) = \mathcal{T}'(P) - 1 < 0$ strictly on $[0, 1]$. By the Banach Contraction Mapping Theorem, $P^*$ is unique and globally stable.
Applying the Implicit Function Theorem to $P^* - \Phi\left(\frac{\tilde{\theta}_t + \lambda \kappa P^*}{\Sigma}\right) = 0$:
\begin{equation}
\frac{\partial P^*}{\partial \tilde{\theta}_t} = \frac{\frac{1}{\Sigma} \phi\left( \frac{\tilde{\theta}_t + \lambda \kappa P^*}{\Sigma} \right)}{1 - \frac{\lambda \kappa}{\Sigma} \phi\left( \frac{\tilde{\theta}_t + \lambda \kappa P^*}{\Sigma} \right)} \equiv \mathcal{M} \cdot \frac{1}{\Sigma}\phi(\cdot)
\end{equation}
Since $\lambda \kappa > 0$, the denominator lies in $(0, 1)$, implying $\mathcal{M} > 1$. Finally, because $\phi(z)$ attains its unique maximum at $z = 0$, $\mathcal{M}$ is maximized when $\tilde{\theta}_t + \lambda \kappa P^* = 0$, which corresponds to $P^* = \Phi(0) = 0.5$. \qed

\subsection{Proof of Proposition \ref{prop:bifurcation}}
\label{app:proof3}
When $\frac{\lambda \kappa}{\Sigma} > \sqrt{2\pi}$, the maximum slope of $\mathcal{T}(P)$ strictly exceeds unity:
\begin{equation}
\max_{P} \mathcal{T}'(P) = \frac{\lambda \kappa}{\sqrt{2\pi}\Sigma} > 1
\end{equation}
Because $\mathcal{T}''(P) > 0$ for $z < 0$ and $\mathcal{T}''(P) < 0$ for $z > 0$, the mapping is strictly sigmoid (S-shaped). For intermediate values of $\tilde{\theta}_t \in (\underline{\theta}, \bar{\theta})$, the 45-degree line intersects $\mathcal{T}(P)$ at exactly three points $0 < P_L^* < P_M^* < P_H^* < 1$. 
At $P_L^*$ and $P_H^*$, the slope satisfies $\mathcal{T}'(P^*) < 1$, establishing local stability. At $P_M^*$, the slope satisfies $\mathcal{T}'(P_M^*) > 1$, establishing instability. Any perturbation past $P_M^*$ induces convergence to the alternate stable branch. \qed

\end{document}
